% ==============================================================================
% BLOQUE 4: GASTÓN GRASSO (Slides 25 a 32 — 8 diapositivas)
% ==============================================================================

% SLIDE 25: SUPERPOSICIÓN DE BODES
\section{Comparativa y Cierre}
\subsection{Superposición de Bodes}
\begin{frame}
  \frametitle{Comparativa de Bodes: Planta Original vs Compensadores}

  \begin{columns}[c, onlytextwidth]
    \begin{column}{0.58\textwidth}
      \centering
      \adjustbox{max width=\linewidth, max height=5.5cm, keepaspectratio}{\begin{tikzpicture}
    \pgfplotsset{
        bode plot/.style={
            /pgf/number format/read comma as period,
            grid=both,
            major grid style={solid, gray!45},
            minor grid style={dotted, gray!35},
            axis line style={black, line width=0.6pt},
            axis x line* = bottom,
            axis y line* = left,
            extra description/.code={
                \draw[gray!45] (rel axis cs:0,1) -- (rel axis cs:1,1) -- (rel axis cs:1,0);
            },
            width=15.5cm,
            height=6.5cm,
            xmin=100, xmax=10000000,
            enlarge x limits=false,
            tick label style={font=\small},
            label style={font=\normalsize},
        }
    }

    % Subplot 1: Magnitud comparada
    \begin{semilogxaxis}[
        bode plot,
        ylabel={Magnitud [\si{\decibel}]},
        ymin=-100, ymax=75,
        ytick={-80, -60, -40, -20, 0, 20, 40, 60},
        xticklabel=\empty,
        legend pos=south west,
        legend style={font=\small, fill=white, fill opacity=0.9, text opacity=1},
        name=magplot,
    ]
        % Planta
        \addplot[line width=1.3pt, blue!85!black] table [col sep=comma, x=frequency, y=mag_planta] {sources/plots/bode_comparativa.csv};
        \addlegendentry{Planta original $G(s)$}

        % Retardo
        \addplot[line width=1.3pt, red!85!black] table [col sep=comma, x=frequency, y=mag_retardo] {sources/plots/bode_comparativa.csv};
        \addlegendentry{Comp. Retardo $G_{\text{ret}}(s)$}

        % Atraso-Adelanto
        \addplot[line width=1.3pt, teal!80!black] table [col sep=comma, x=frequency, y=mag_aa] {sources/plots/bode_comparativa.csv};
        \addlegendentry{Comp. Atraso-Adelanto $G_{\text{aa}}(s)$}

        % Nivel 0 dB
        \addplot[black!70, densely dotted, line width=0.8pt] coordinates {(100, 0) (10000000, 0)};
        \node[black!70, anchor=south east, font=\scriptsize] at (axis cs:9000000, 1.5) {$0\,\text{dB}$};

        % Proyecciones verticales f_45
        \draw[blue!85!black, dashed, line width=0.9pt] (axis cs:52901, -100) -- (axis cs:52901, 43.44);
        \draw[red!85!black, dashed, line width=0.9pt] (axis cs:9962, -100) -- (axis cs:9962, 38.25);
        \draw[teal!80!black, dashed, line width=0.9pt] (axis cs:45109, -100) -- (axis cs:45109, 45.18);

        % Marcadores f_45
        \addplot[only marks, mark=*, mark size=2.5pt, blue!85!black] coordinates {(52901, 43.44)};
        \node[blue!85!black, fill=white, fill opacity=0.9, text opacity=1, inner sep=1.5pt, font=\scriptsize, anchor=south west] 
            at (axis cs:56000, 44.5) {$43,44\,\text{dB}$};

        \addplot[only marks, mark=*, mark size=2.5pt, red!85!black] coordinates {(9962, 38.25)};
        \node[red!85!black, fill=white, fill opacity=0.9, text opacity=1, inner sep=1.5pt, font=\scriptsize, anchor=north] 
            at (axis cs:9962, 36) {$38,25\,\text{dB}$};

        \addplot[only marks, mark=*, mark size=2.5pt, teal!80!black] coordinates {(45109, 45.18)};
        \node[teal!80!black, fill=white, fill opacity=0.9, text opacity=1, inner sep=1.5pt, font=\scriptsize, anchor=south east] 
            at (axis cs:42000, 46.5) {$45,18\,\text{dB}$};
    \end{semilogxaxis}

    % Subplot 2: Fase comparada (posicionado estrictamente DEBAJO de magplot)
    \begin{semilogxaxis}[
        bode plot,
        xlabel={$f$ [\si{\hertz}]},
        ylabel={Fase [$^\circ$]},
        ymin=-370, ymax=10,
        ytick={-360, -315, -270, -225, -180, -135, -90, -45, 0},
        at={(magplot.south west)},
        anchor=north west,
        yshift=-0.7cm,
        legend pos=south west,
        legend style={font=\small, fill=white, fill opacity=0.9, text opacity=1},
    ]
        % Planta
        \addplot[line width=1.3pt, blue!85!black] table [col sep=comma, x=frequency, y=phase_planta] {sources/plots/bode_comparativa.csv};
        \addlegendentry{Planta original}

        % Retardo
        \addplot[line width=1.3pt, red!85!black] table [col sep=comma, x=frequency, y=phase_retardo] {sources/plots/bode_comparativa.csv};
        \addlegendentry{Retardo}

        % Atraso-Adelanto
        \addplot[line width=1.3pt, teal!80!black] table [col sep=comma, x=frequency, y=phase_aa] {sources/plots/bode_comparativa.csv};
        \addlegendentry{Atraso-Adelanto}

        % Referencia -135°
        \addplot[purple, dashed, line width=1.0pt] coordinates {(100, -135) (10000000, -135)};
        \node[purple, fill=white, fill opacity=0.85, text opacity=1, anchor=south east, font=\small, inner sep=1.5pt] 
            at (axis cs:9000000, -132) {$-135^\circ\;(MF=45^\circ)$};

        % Referencia -180°
        \addplot[red!80!black, dotted, line width=1.0pt] coordinates {(100, -180) (10000000, -180)};
        \node[red!80!black, fill=white, fill opacity=0.85, text opacity=1, anchor=south east, font=\small, inner sep=1.5pt] 
            at (axis cs:9000000, -177) {$-180^\circ$};

        % Proyecciones verticales f_45
        \draw[blue!85!black, dashed, line width=0.9pt] (axis cs:52901, -370) -- (axis cs:52901, -135);
        \draw[red!85!black, dashed, line width=0.9pt] (axis cs:9962, -370) -- (axis cs:9962, -135);
        \draw[teal!80!black, dashed, line width=0.9pt] (axis cs:45109, -370) -- (axis cs:45109, -135);

        % Cruces en f_45
        \addplot[only marks, mark=*, mark size=2.5pt, blue!85!black] coordinates {(52901, -135)};
        \node[blue!85!black, fill=white, fill opacity=0.85, text opacity=1, inner sep=1pt, font=\scriptsize, anchor=south west] 
            at (axis cs:56000, -130) {$f_{45}=52,90\,\text{kHz}$};

        \addplot[only marks, mark=*, mark size=2.5pt, red!85!black] coordinates {(9962, -135)};
        \node[red!85!black, fill=white, fill opacity=0.85, text opacity=1, inner sep=1pt, font=\scriptsize, anchor=north] 
            at (axis cs:9962, -143) {$f_{45}=9,96\,\text{kHz}$};

        \addplot[only marks, mark=*, mark size=2.5pt, teal!80!black] coordinates {(45109, -135)};
        \node[teal!80!black, fill=white, fill opacity=0.85, text opacity=1, inner sep=1pt, font=\scriptsize, anchor=north east] 
            at (axis cs:42000, -143) {$f_{45}=45,11\,\text{kHz}$};
    \end{semilogxaxis}
\end{tikzpicture}
}
    \end{column}

    \begin{column}{0.40\textwidth}
      \begin{tarjeta}
        {\fontsize{8.5pt}{10pt}\selectfont \textbf{\color{NavyDark} Análisis Visual Comparativo:}}\\[2pt]
        {\fontsize{7pt}{8.5pt}\selectfont
        \begin{itemize}\setlength{\itemsep}{2pt}
          \item \textbf{Curva Azul:} Planta original ($BW = \qty{11.59}{\kilo\hertz}$).
          \item \textbf{Curva Roja:} Retardo. Caída prematura desde $1.2\text{ kHz}$; contrae el ancho de banda a $\qty{1.19}{\kilo\hertz}$ (pérdida del $\qty{89.7}{\percent}$).
          \item \textbf{Curva Verde:} Atraso-Adelanto. Preserva la respuesta útil ($BW = \qty{11.56}{\kilo\hertz}$) y atenúa únicamente en alta frecuencia.
        \end{itemize}
        }
      \end{tarjeta}
    \end{column}
  \end{columns}
\end{frame}

% SLIDE 26: CUADRO COMPARATIVO CUANTITATIVO
\subsection{Cuadro Comparativo}
\begin{frame}
  \frametitle{Cuadro Comparativo Cuantitativo Integral (MF = 45°)}

  \begin{tarjeta}
    \centering
    \setlength{\tabcolsep}{3.2pt}
    \renewcommand{\arraystretch}{1.08}
    {\fontsize{7.2pt}{9.0pt}\selectfont
    \begin{tabular}{lccc}
      \toprule
      \textbf{Especificación / Parámetro} & \textbf{Planta Original} & \textbf{Comp. Retardo} & \textbf{Comp. Atraso-Adelanto} \\
      \midrule
      Ganancia estática continua ($K$) & $\qty{59.08}{\decibel}$ & $\qty{59.08}{\decibel}$ & $\qty{59.08}{\decibel}$ \\
      Ancho de banda a $\qty{-3}{\decibel}$ ($BW_{\text{LA}}$) & $\qty{11.59}{\kilo\hertz}$ & $\qty{1.19}{\kilo\hertz}$ & $\mathbf{\qty{11.56}{\kilo\hertz}}$ \\
      Variación porcentual ancho de banda & Base & \textbf{-89.7\%} & \textbf{-0.3\%} \\
      Frecuencia para $MF = \ang{45}$ ($f_{45}$) & $\qty{52.90}{\kilo\hertz}$ & $\qty{9.96}{\kilo\hertz}$ & $\mathbf{\qty{45.11}{\kilo\hertz}}$ \\
      Ganancia de lazo disponible en $f_{45}$ & $\qty{43.44}{\decibel}$ & $\qty{38.25}{\decibel}$ & $\qty{45.18}{\decibel}$ \\
      Frecuencia de cruce por $\ang{-180}$ ($f_{180}$) & $\qty{110.99}{\kilo\hertz}$ & $\qty{24.35}{\kilo\hertz}$ & $\qty{84.49}{\kilo\hertz}$ \\
      Ganancia en cruce de fase $\ang{-180}$ & $\qty{32.14}{\decibel}$ & $\qty{25.21}{\decibel}$ & $\qty{36.22}{\decibel}$ \\
      Margen de Ganancia disponible ($MG_{45}$) & $\qty{11.30}{\decibel}$ & $\mathbf{\qty{13.04}{\decibel}}$ & $\mathbf{\qty{8.96}{\decibel}}$ \\
      Margen de Ganancia directo sin atenuar & $\qty{-32.14}{\decibel}$ & $\qty{-25.21}{\decibel}$ & $\qty{-36.22}{\decibel}$ \\
      \bottomrule
    \end{tabular}
    }
  \end{tarjeta}
\end{frame}

% SLIDE 27: IMPACTO EN ANCHO DE BANDA
\subsection{Análisis Comparativo}
\begin{frame}
  \frametitle{Impacto sobre el Ancho de Banda y Rapidez de Respuesta}

  \begin{columns}[t, onlytextwidth]
    \begin{column}{0.48\textwidth}
      \begin{tarjeta}
        {\fontsize{9pt}{11pt}\selectfont \textbf{\color{NavyDark} Compensador de Retardo (Lag)}}\\[2pt]
        {\fontsize{8pt}{10pt}\selectfont
        \begin{itemize}\setlength{\itemsep}{2.5pt}
          \item \textbf{Ancho de banda base de la planta:}
          $BW = \qty{11.59}{\kilo\hertz}$, gobernado por su primer polo ($f_1 = \qty{12}{\kilo\hertz}$).
          \item \textbf{Acción del polo dominante:}
          Introduce polo en $f_d = \qty{1.2}{\kilo\hertz}$, forzando caída prematura de $-20\text{ dB/década}$.
          \item \textbf{Contracción drástica de ancho de banda:}
          $BW = f_{-3\text{dB}} \approx \mathbf{\qty{1.19}{\kilo\hertz}}$ (pérdida del $\mathbf{89.7\%}$).
          \item \textbf{Frecuencia para $MF = \ang{45}$:}
          Se alcanza a baja velocidad: $f_{45} \approx \mathbf{\qty{9.96}{\kilo\hertz}}$.
        \end{itemize}
        }
      \end{tarjeta}
    \end{column}

    \begin{column}{0.48\textwidth}
      \begin{tarjeta}[colframe=BlueAccent]
        {\fontsize{9pt}{11pt}\selectfont \textbf{\color{BlueAccent} Compensador Atraso-Adelanto}}\\[2pt]
        {\fontsize{8pt}{10pt}\selectfont
        \begin{itemize}\setlength{\itemsep}{2.5pt}
          \item \textbf{Frecuencias de la red pasiva:}
          Polo $f_p = \qty{220}{\kilo\hertz}$ y cero $f_z = \qty{1.1}{\mega\hertz}$, ambos por encima de la planta.
          \item \textbf{Preservación de banda base:}
          La ganancia y fase en bajas frecuencias se mantienen prácticamente intactas.
          \item \textbf{Retención del espectro útil:}
          $BW = f_{-3\text{dB}} \approx \mathbf{\qty{11.56}{\kilo\hertz}}$ (retiene el $\mathbf{99.7\%}$).
          \item \textbf{Frecuencia para $MF = \ang{45}$:}
          $f_{45} \approx \mathbf{\qty{45.11}{\kilo\hertz}}$ ($4.5$ veces más veloz).
        \end{itemize}
        }
      \end{tarjeta}
    \end{column}
  \end{columns}
\end{frame}

% SLIDE 28: MÁRGENES DE ESTABILIDAD
\subsection{Márgenes de Estabilidad}
\begin{frame}
  \frametitle{Margen de Ganancia y Atenuación en Alta Frecuencia}

  \begin{columns}[t, onlytextwidth]
    \begin{column}{0.48\textwidth}
      \begin{tarjeta}
        {\fontsize{9pt}{11pt}\selectfont \textbf{\color{NavyDark} Márgenes Obtenidos para $MF = \ang{45}$}}\\[2pt]
        {\fontsize{8pt}{10pt}\selectfont
        \begin{itemize}\setlength{\itemsep}{2.5pt}
          \item \textbf{Planta sin compensar:}
          Margen de ganancia directo fuertemente negativo ($MG = \mathbf{-32.14\ \unit{\decibel}}$).
          \item \textbf{Compensador de retardo:}
          Margen disponible $MG_{45} = \mathbf{\qty{13.04}{\decibel}}$ (cruce por $\ang{-180}$ en $f_{180} = \qty{24.35}{\kilo\hertz}$).
          \item \textbf{Compensador atraso-adelanto:}
          Margen disponible $MG_{45} = \mathbf{\qty{8.96}{\decibel}}$ (cruce por $\ang{-180}$ en $f_{180} = \qty{84.49}{\kilo\hertz}$).
          \item \textbf{Criterio estándar de ingeniería:}
          Ambas redes verifican holgadamente la condición de estabilidad robusta: $MG \ge \qty{6}{\decibel}$.
        \end{itemize}
        }
      \end{tarjeta}
    \end{column}

    \begin{column}{0.48\textwidth}
      \begin{tarjeta}
        {\fontsize{9pt}{11pt}\selectfont \textbf{\color{NavyDark} Atenuación y Dinámica de Fase}}\\[2pt]
        {\fontsize{8pt}{10pt}\selectfont
        \begin{itemize}\setlength{\itemsep}{2.5pt}
          \item \textbf{Retardo RC:}
          Mantiene pendiente de $-20\text{ dB/década}$ en todo el rango superior; máxima atenuación de alta frecuencia.
          \item \textbf{Atraso-adelanto:}
          La atenuación está acotada entre $f_p$ y $f_z$ por $\frac{R_2}{R_1+R_2} = \frac{1}{5}$ ($\qty{-13.98}{\decibel}$).
          \item \textbf{Aporte del cero ($1.1\text{ MHz}$):}
          Frena la caída de magnitud y recupera fase, estabilizando el lazo sin sacrificar la velocidad base.
          \item \textbf{Compromiso:} Retardo entrega $\qty{4.08}{\decibel}$ más de margen, pero atraso-adelanto cuadruplica la velocidad.
        \end{itemize}
        }
      \end{tarjeta}
    \end{column}
  \end{columns}
\end{frame}

% SLIDE 29: VENTAJAS POR RÉGIMEN
\subsection{Regímenes de Amortiguamiento}
\begin{frame}
  \frametitle{Ventajas y Desventajas por Régimen de Amortiguamiento}

  \begin{columns}[t, onlytextwidth]
    \begin{column}{0.48\textwidth}
      \begin{tarjeta}
        {\fontsize{8.5pt}{10.5pt}\selectfont \textbf{\color{NavyDark} Sobreamortiguado ($\xi = 2.0$, $Q = 0.25$)}}\\[2pt]
        {\fontsize{7.5pt}{9.5pt}\selectfont
        \begin{itemize}\setlength{\itemsep}{2pt}
          \item \textbf{Sin sobreimpulso:} Dinámica monotónica sin \emph{ringing} (subpico despreciable de $\mathbf{-4.8\%}$).
          \item \textbf{Alta robustez:} Muy tolerante frente a derivas y dispersión en $L$, $C$ o $R_g$.
          \item \textbf{Compromiso de lentitud:} Establecimiento prolongado ($t_s = \mathbf{\qty{7.53}{\milli\second}}$, $\tau_1 = \qty{5.57}{\milli\second}$).
          \item \textbf{Baja selectividad:} Transición suave de $-20\text{ dB/década}$ hasta el segundo polo.
          \item \textbf{Componentes:} Exige capacidad elevada ($C = \qty{59.70}{\micro\farad}$).
        \end{itemize}
        }
      \end{tarjeta}
    \end{column}

    \begin{column}{0.48\textwidth}
      \begin{tarjeta}
        {\fontsize{8.5pt}{10.5pt}\selectfont \textbf{\color{NavyDark} Subamortiguado ($\xi = 0.667$, $Q = 0.75$)}}\\[2pt]
        {\fontsize{7.5pt}{9.5pt}\selectfont
        \begin{itemize}\setlength{\itemsep}{2pt}
          \item \textbf{Elevada rapidez:} Caída rápida ($t_{\text{caída}} = \qty{0.33}{\milli\second}$) y establecimiento $4\times$ menor ($t_s = \mathbf{\qty{1.82}{\milli\second}}$).
          \item \textbf{Alta selectividad:} Caída abrupta de $-40\text{ dB/década}$ inmediatamente bajo $f_c$.
          \item \textbf{Componentes compactos:} Capacidad $12\times$ menor ($C = \qty{5.02}{\micro\farad}$).
          \item \textbf{Subpico pronunciado:} Oscilación amortiguada con subpico del $\mathbf{-22.2\%}$ ($T_d = \qty{3.17}{\milli\second}$).
          \item \textbf{Crítico ($Q = 0.5$):} Óptimo sin sobreimpulso pero sensible a tolerancias.
        \end{itemize}
        }
      \end{tarjeta}
    \end{column}
  \end{columns}
\end{frame}

% SLIDE 30: CRITERIOS DE ELECCIÓN
\subsection{Aplicación en Ingeniería}
\begin{frame}
  \frametitle{Criterios de Elección según la Aplicación}

  \begin{columns}[t, onlytextwidth]
    \begin{column}{0.48\textwidth}
      \begin{tarjeta}
        {\fontsize{9pt}{11pt}\selectfont \textbf{\color{NavyDark} Régimen de Amortiguamiento}}\\[2pt]
        {\fontsize{8pt}{10pt}\selectfont
        \begin{itemize}\setlength{\itemsep}{2.5pt}
          \item \textbf{Sobreamortiguado:}
          Recomendado cuando la señal alimenta etapas sensibles a excursiones de signo contrario (comparadores, detectores de nivel, limitadores de excursión), donde un sobreimpulso generaría disparos espurios.
          \item \textbf{Subamortiguado:}
          Óptimo para canales de transmisión selectivos y procesamiento de pulsos rápidos donde se priorice un tiempo de establecimiento mínimo ($t_s < \qty{2}{\milli\second}$) y atenuación abrupta en rechazo.
        \end{itemize}
        }
      \end{tarjeta}
    \end{column}

    \begin{column}{0.48\textwidth}
      \begin{tarjeta}
        {\fontsize{9pt}{11pt}\selectfont \textbf{\color{NavyDark} Red de Compensación}}\\[2pt]
        {\fontsize{8pt}{10pt}\selectfont
        \begin{itemize}\setlength{\itemsep}{2.5pt}
          \item \textbf{Red de Atraso-Adelanto:}
          Óptima en etapas de amplificación de instrumentación o procesamiento de señales donde el tiempo de establecimiento sea un parámetro crítico, requiriendo retener el $\qty{99.7}{\percent}$ del ancho de banda útil con $MG \ge \qty{6}{\decibel}$.
          \item \textbf{Compensador de Retardo RC:}
          Alternativa idónea en aplicaciones donde prime la inmunidad al ruido y la tolerancia frente a derivas térmicas o dispersión de fabricación, a costa de una respuesta temporal perceptiblemente más lenta.
        \end{itemize}
        }
      \end{tarjeta}
    \end{column}
  \end{columns}
\end{frame}

% SLIDE 31: CONCLUSIONES
\subsection{Conclusiones Integrales}
\begin{frame}
  \frametitle{Conclusiones Integrales del Proyecto}

  \begin{columns}[t, onlytextwidth]
    \begin{column}{0.48\textwidth}
      \begin{tarjeta}
        {\fontsize{9pt}{11pt}\selectfont \textbf{\color{NavyDark} Filtro Pasa Alto de Segundo Orden}}\\[2pt]
        {\fontsize{8pt}{10pt}\selectfont
        \begin{itemize}\setlength{\itemsep}{2.5pt}
          \item \textbf{Demostración de $\omega_0 \neq \omega_c$:} La constante de corrección $K$ fue indispensable para validar analíticamente los $\qty{-3.01}{\decibel}$ en $\qty{400}{\hertz}$.
          \item \textbf{Correspondencia unívoca:} Las raíces en plano $s$ describen rigurosamente la dinámica observada en el dominio del tiempo.
          \item \textbf{Compromiso resuelto:} Amortiguamiento elevado brinda suavidad y robustez; amortiguamiento bajo ofrece rapidez ($4\times$) y selectividad.
        \end{itemize}
        }
      \end{tarjeta}
    \end{column}

    \begin{column}{0.48\textwidth}
      \begin{tarjeta}
        {\fontsize{9pt}{11pt}\selectfont \textbf{\color{NavyDark} Control y Estabilidad de Tercer Orden}}\\[2pt]
        {\fontsize{8pt}{10pt}\selectfont
        \begin{itemize}\setlength{\itemsep}{2.5pt}
          \item \textbf{Inestabilidad de planta pura:} Ganancia de $K=900$ con tres polos reales arroja $MG = \qty{-32.14}{\decibel}$, demandando compensación obligatoria.
          \item \textbf{Criterio $MF = \ang{45}$ verificado:} Ambas redes garantizan estabilidad robusta ($MG \ge \qty{6}{\decibel}$) y sobreimpulso acotado a $< \qty{25}{\percent}$.
          \item \textbf{Superioridad de atraso-adelanto:} Preserva el $\qty{99.7}{\percent}$ del ancho de banda base ($11.56\text{ kHz}$) y opera a $f_{45} = \qty{45.11}{\kilo\hertz}$.
        \end{itemize}
        }
      \end{tarjeta}
    \end{column}
  \end{columns}
\end{frame}
